\documentclass[12pt]{report}
\usepackage{amsmath}
\usepackage{amssymb}
\usepackage{amsthm}
\usepackage{enumitem}
\usepackage{multicol}

\newenvironment{case}[2]{\newline\textbf{Case #1: #2.}}{}

\begin{document}
\renewcommand\qedsymbol{OE$\Delta$}
\newcommand{\abs}[1]{\left| #1 \right|}
\newcommand{\andtxt}{\text{ and }}
\newcommand{\mathif}{\text{if }}
\newcommand{\seq}[2][n]{\ensuremath{(#2 _{#1})}}

\tableofcontents
\pagebreak

\setcounter{chapter}{3}
\chapter{Series not on TV}

\section{Problem 4.1}
    Determine whether each of the following converges conditionally, converges absolutely, or diverges. 
    You do not need  to prove your answers, but state which of the following tests gives the answers: the $k^{\rm th}$-term test, the geometric series test, and the alternating series test.
    \begin{multicols}{3}
        \begin{enumerate}[label=(\alph*)]
            \item   $\sum_{k = 1}^{\infty}  (-1)^k \frac{1}{\sqrt{k}}$
            \item   $\sum_{k = 1}^{\infty}  (-1)^k \frac{k}{k + 7}$
            \item   $\sum_{k = 1}^{\infty}  \frac{1}{(\ln(4))^k}$
            \item   $\sum_{k = 1}^{\infty}  \frac{1}{k!}$
            \item   {\scriptsize $\sum_{k = 1}^{\infty}  (\sqrt{k + 1} - \sqrt{k})$}
            \item   $\sum_{k = 1}^{\infty}  \frac{1}{2}$
            \item   $\sum_{k = 1}^{\infty}  (-1)^k \frac{1}{\sqrt[3]{k}}$
            \item   $\sum_{k = 53}^{\infty} \frac{\ln(k)}{\ln(\ln(k))}$
            \item   $\sum_{k = 7}^{\infty}  \left( \frac{2}{k} + \frac{3}{k^2} \right)$
        \end{enumerate}
    \end{multicols}

\section{Problem 4.3}
    Prove that if $\sum_{k = 1}^{\infty} \abs{a_k}$ converges, then $\sum_{k = 1}^{\infty} a_k$ converge too.
    This is Propostion 4.18.

\section{Problem 4.5}
    Give an example of a series where $\sum_{k = 1}^{\infty} a_k$ converges but $\sum_{k = 1}^{\infty} a_{2k}$ diverges.

\section{Problem 4.7}
    Give an example of a series $\sum_{k = 1}^{\infty} a_k$ where:
    \begin{multicols}{3}
        \begin{itemize}
            \item {\small $\sum_{k = 1}^{\infty} a_k$ converges,}
            \item {\small $\sum_{k = 1}^{\infty} a_k^2$ diverges}
            \item $\sum_{k = 1}^{\infty} a_k^3$ converges
        \end{itemize}
    \end{multicols}

\section{Problem 4.9}
    The geometric series test (Proposition 4.9) says  that a geometric series diverges if $\abs{r} \geq 1$. 
    Recall that a series can  diverge to $\infty$, to $-\infty$, or can be ``does not exist.'' 
    Which form of divergence is it for a geometric series? 
    Note that your answer may depend on r. 

\section{Problem 4.11}
    \begin{enumerate}[label=(\alph*)]
        \item   Find a way to write 77.77777777... as a geometric series, and then prove this number is rational by using the geometric series test to write this number as a fraction with integers in the numerator and denominator.
        \item   Write 77.77777777... as a different geometric series, and use the geometric series test to write this number as a fraction with integers in the numerator and denominator. Are your two fractions the same?
        \item   A number $q$ has a repeating decimal if the non-integer portion of its decimal expansion is repetitive. For example, 72.578578578578578... has a repeating decimal. Prove that if a number $q$ has a repeating  decimal, then $q$ is rational. 
    \end{enumerate}
\pagebreak

\section{Problem 4.13}
    \begin{enumerate}[label=(\alph*)]
        \item   Prove that if $(k a_k)$ converges to a nonzero real number $L$, then  the series $\sum_{k = 1}^{\infty} a_k$ diverges. Give an example to show that the  converse is false. 
        \item   Prove that if $(k^2 a_k)$ converges (to any real number), then the  series $\sum_{k = 1}^{\infty} a_k$ converges. Give an example to show that the  converse is false. 
    \end{enumerate}

    \subsection{Solution (a)}
        \begin{proof}
            We know that \seq[k]{ka} converges to a nonzero $L$.
            \begin{equation}
                \forall\varepsilon > 0, \exists N : \forall k > N, \abs{ka_k - L} < \varepsilon
            \end{equation}
    
            It is very much possible that $\varepsilon = \frac{\abs{L}}{2}$.
            We can also apply the reverse triangle inequality and isolate $\abs{a_k}$.
            \begin{gather}
                \frac{\abs{L}}{2} = \varepsilon > \abs{ka_k - L} \geq \abs{\abs{ka_k} - \abs{L}}\\
                \frac{\abs{L}}{2} > \abs{ka_k} - \abs{L} > -\frac{\abs{L}}{2}\\
                \frac{3\abs{L}}{2} > \abs{ka_k} > \frac{\abs{L}}{2}\\
                \frac{3\abs{L}}{2} > k\abs{a_k} > \frac{\abs{L}}{2}\\
                \frac{3\abs{L}}{2k} > \abs{a_k} > \frac{\abs{L}}{2k} = \frac{\abs{L}}{2} \cdot \frac{1}{k}
            \end{gather}
    
            Because of its convergence to $L$, beyond a certain point, $a_k$ has the same sign as $L$.
            If we only work beyond that point, we can remove the absolute values.
            This gives us two cases.
            First is the case where $L > 0$.
            \begin{gather}
                \frac{3L}{2k} > a_k > \frac{L}{2} \cdot \frac{1}{k}
            \end{gather}
    
            The other case is the one where $L < 0$.
            \begin{gather}
                -\frac{3L}{2k} > -a_k > -\frac{L}{2} \cdot \frac{1}{k}
            \end{gather}
    
            We can substitute in a new sequence \seq[k]{b} in case 2 where for any $k$, $b_k = -a_k$ to get back to the point described in case 1.
            Since it's the same point, I won't be replacing $a_k$ with $b_k$.
            \begin{gather}
                \frac{3L}{2k} > a_k > \frac{L}{2} \cdot \frac{1}{k}
            \end{gather}
    
            Now, recall the harmonic series.
            \begin{equation}
                \sum_{k=1}^{\infty} \frac{1}{k} \to \infty
            \end{equation}
    
            We can multiply this in both cases by $\frac{L}{2}$.
            \begin{equation}
                \sum_{k=1}^{\infty} \frac{L}{2} \cdot \frac{1}{k} \to \infty
            \end{equation}
            
            Since the harmonic series diverges and \seq[k]{a} is always greater than our above alteration of the harmonic series, $\sum_{k = 1}^{\infty} a_k$ diverges.
        \end{proof}

        For an example of how the converse could be false, consider the series $\sum_{k = 1}^{\infty} k$.
        It diverges, and the components of it, even when multiplied by $k$, also diverge.
    
    \subsection{Solution (b)}
        \begin{proof}
            We know that \seq[k]{k^2 a} converges to a real number $L$.
            \begin{equation}
                \forall\varepsilon > 0, \exists N : \forall k > N, \abs{k^2 a_k - L} < \varepsilon
            \end{equation}
    
            We consider this based on whather $L$ is or is not equal to zero.
            \begin{case}{1}{$L \neq 0$}
                If $L \neq 0$, it is very much possible that $\varepsilon = \frac{\abs{L}}{2}$.
                We can also apply the reverse triangle inequality and isolate $\abs{a_k}$.
                \begin{gather}
                    \frac{\abs{L}}{2} = \varepsilon > \abs{k^2 a_k - L} \geq \abs{\abs{k^2 a_k} - \abs{L}}\\
                    \frac{\abs{L}}{2} > \abs{k^2 a_k} - \abs{L} > -\frac{\abs{L}}{2}\\
                    \frac{3\abs{L}}{2} > \abs{k^2 a_k} > \frac{\abs{L}}{2}\\
                    \frac{3\abs{L}}{2} > k^2\abs{a_k} > \frac{\abs{L}}{2}\\
                    \frac{3\abs{L}}{2k^2} = \frac{3\abs{L}}{2} \cdot \frac{1}{k^2} > \abs{a_k} > \frac{\abs{L}}{2k^2} = \frac{\abs{L}}{2} \cdot \frac{1}{k^2}
                \end{gather}

                By the series p-test, we can recognize the convergence of $\sum_{k = 1}^{\infty} \frac{1}{k^2}$, and multiply that by a constant.
                \begin{gather}
                    \sum_{k = 1}^{\infty} \frac{1}{k^2} \to 0\\
                    \sum_{k = 1}^{\infty} \frac{3\abs{L}}{2} \cdot \frac{1}{k^2} \to 0
                \end{gather}

                By the comparison test, this means that $\sum_{k=1}^{\infty} \abs{a_k}$ converges.
                This in turn means that $\sum_{k=1}^{\infty} a_k$ converges.
            \end{case}

            Now for the other case. 
            \begin{case}{2}{$L = 0$}
                We can rewrite the initial definition of convergence based on this.
                \begin{equation}
                    \forall\varepsilon > 0, \exists N : \forall k > N, \abs{k^2 a_k} < \varepsilon
                \end{equation}

                % Recall that $k > 0$, and by extension that $k^2 > 0$.
                % \begin{gather}
                %     k^2 > 0\\
                %     k^2 \abs{a_k} > \abs{a_k}
                % \end{gather}

                % This means that \seq[]{\abs{a_k}} converges to zero.
                Now, suppose that $\varepsilon = 1$.
                Solve for $\abs{a_k}$.
                \begin{gather}
                    \abs{k^2 a_k} < 1\\
                    \abs{a_k} < \frac{1}{k^2}
                \end{gather}

                As with earlier, this means the sum converges.
            \end{case}
        \end{proof}

        For the example disproving the converse, consider \seq[]{\sum_{k=1}^{\infty}k^{-\frac{3}{2}}}.
        This converges.
        However, $(k^2 a_k)$ is equal to $(\sqrt{k})$, which diverges to infinity.
    \pagebreak

\section{Problem 4.15}
    Prove the Cauchy condensation test. 
    That is, suppose that \seq{a} is a decreasing sequence for which $a_k \to 0$. 
    Prove that $\sum_{k=1}^{\infty} a_k$ converges if and only if $\sum_{k=1}^{\infty} 2^k a_{2^k}$ converges. 

    \subsection{Solution}
        First, prove that if $\sum_{k=1}^{\infty} a_k$ converges, $\sum_{k=1}^{\infty} 2^k a_{2^k}$ converges.
        \begin{proof}
            Start with convergent $\sum_{k=1}^{\infty} a_k$ for decreasing $a_k \to 0$.
            We can recall the definition of the two.
            \begin{gather}
                \forall \varepsilon > 0, \exists N : \forall k > N, \abs{a_k} < \varepsilon\\
                a_{k + 1} < a_{k}\\
                \sum_{k=1}^{\infty} a_k \to L_1
            \end{gather}

            Here, we can try grouping terms in \seq[k]{a}.
            There is an inequality that is extrapolated from this being a decreasing sequence.
            \begin{gather}
                \forall n \in \mathbb{N}, 2^{n+1} - 1 \geq j \geq 2^n \implies a_{2^n} \geq a_j \geq a_{2^{n+1} - 1}\\
                \sum_{k=2^n}^{2^{n+1} - 1} a_k \geq 2^{n} a_{2^{n + 1}} \geq 0
            \end{gather}

            We can furthermore sum this over $n$.
            \begin{gather}
                \sum_{n = 1}^{\infty} \sum_{k=2^n}^{2^{n+1} - 1} a_k \geq \sum_{n = 1}^{\infty} 2^{n} a_{2^{n + 1}} \geq 0
            \end{gather}

            The former term is equal to the converging $\sum_{k=1}^{\infty} a_k$.
            Therefore, the latter term also converges.
            The latter term can be turned into $\sum_{k=1}^{\infty} 2^k a_{2^k}$ by adding a few terms.
            \begin{equation}
                \sum_{k=1}^{\infty} 2^k a_{2^k} = 2*\sum_{n = 1}^{\infty} 2^{n} a_{2^{n + 1}} + 2 * a_2
                    =   \sum_{n = 1}^{\infty} 2^{n + 1} a_{2^{n + 1}} + 2 * a_2
            \end{equation}

            None of these operations (multiplying by 2 and adding $2 * a_2$) compromise convergence.
            Ergo, $\sum_{k=1}^{\infty} 2^k a_{2^k}$ converges.
        \end{proof}

        Next, prove that if $\sum_{k=1}^{\infty} 2^k a_{2^k}$ converges, then $\sum_{k=1}^{\infty} a_k$ converges.
        \begin{proof}
            We begin with the earlier statement.
            \begin{gather}
                \forall n \in \mathbb{N}, 2^{n+1} - 1 \geq j \geq 2^n \implies a_{2^n} \geq a_j \geq a_{2^{n+1} - 1}\\
                2^{n} a_{2^{n}} \geq \sum_{k=2^n}^{2^{n+1} - 1} a_k \geq 0
            \end{gather}

            Sum this over $n$.
            \begin{gather}
                \sum_{n=1}^{\infty} 2^{n} a_{2^{n}} \geq \sum_{n=1}^{\infty} \sum_{k=2^n}^{2^{n+1} - 1} a_k \geq 0
            \end{gather}

            The latter term converges. 
            As described earlier, it is equal to the sum over $\mathbb{N}$ of \seq[k]{a}.
            \begin{equation}
                \sum_{n=1}^{\infty} \sum_{k=2^n}^{2^{n+1} - 1} a_k = \sum_{k=1}^{\infty} a_k
            \end{equation}

            Therefore, $\sum_{k=1}^{\infty} a_k$ converges.
        \end{proof}
    \pagebreak

\section{Problem 4.17}
    Prove the ratio test via the following steps. 
    Given a series $\sum_{k = 1}^{\infty} a_k$ with $a_k = 0$, assume that 
    \begin{equation*}
        \lim_{k \to \infty} \abs{\frac{a_{k+1}}{a_{k}}} =: r < 1
    \end{equation*}
    We will prove that the series converges absolutely.
    \begin{enumerate}[label=(\alph*)]
        \item   Let $q$ be such that $r < q < 1$. Explan why there is some $N$ such that $n \geq N$ implies that $\abs a_{k+1} \leq \abs{a_k} \cdot q$.
        \item   Explain why $\sum_{k = 1}^{\infty} \abs{a_N} \cdot q^k$ necessarily converges.
        \item   Finally, use part (b) to prove that $\sum_{k = 1}^{\infty} \abs{a_k}$ converges.
    \end{enumerate}
    \pagebreak

\section{Problem 4.18}
    Prove the root test via the following steps. 
    Given a series $\sum_{k=1}^{\infty} a_k$ where each $a_k \geq 0$, assume that the limit $\lim_{k \to \infty}(a_k)^{1/k}$ exists. 
    Call this limit $\rho$. 
    Then this series converges if $\rho < 1$ and  diverges if $\rho > 1$. 
    (The test is inconclusive if $\rho = 1$.)  
    \begin{enumerate}[label=(\alph*)]
        \item   Suppose $\rho < 1$. Let $\varepsilon = \frac{1 - \rho}{2}$ and $\rho_1 = \rho + \varepsilon$. Prove that there is some $N$ for which $(a_n)^{1/n} < \rho_1$ for all $n>N$.
        \item   Prove that $\sum_{k = N}^{\infty} a_k$ converges by comparing it to a geometric series. Then conclude that $\sum_{k = 1}^{\infty} a_k$ also converges. 
        \item   Suppose $\rho > 1$. Let $\varepsilon = \frac{\rho - 1}{2}$ and $\rho_2 = \rho - \varepsilon$. Prove that there is some $N$ for which $(a_n)^{1/n} > \rho_2$ for all $n>N$.
        \item   Prove that $\sum_{k = N}^{\infty} a_k$ diverges by comparing it to a geometric series. Then conclude that $\sum_{k = 1}^{\infty} a_k$ also diverges.
    \end{enumerate}

    \subsection{Solution (a)}
        Begin with the limit from earlier and recall its convergence definition (traditional and Cauchy).
        \begin{gather}
            \lim_{k \to \infty}(a_k)^{1/k} = \rho < 1\\
            \forall\varepsilon > 0, \exists N: \forall n > N, \abs{(a_n)^{1/n} - \rho} < \varepsilon\\
            \forall\varepsilon > 0, \exists N: \forall m, n > N, \abs{(a_m)^{1/m} - (a_n)^{1/n}} < \varepsilon
        \end{gather}

        Expand the first absolute value, then isolate $(a_n)^{1/n}$.
        \begin{gather}
            -\varepsilon < (a_n)^{1/n} - \rho < \varepsilon\\
            \rho - \varepsilon < (a_n)^{1/n} < \rho + \varepsilon
        \end{gather}
    \pagebreak

\section{Problem 4.20}
    \begin{enumerate}[label=(\alph*)]
        \item   Give an example of two divergent series $\sum_{k=1}^{\infty} a_k$ and $\sum_{k=1}^{\infty} b_k$, such that $\sum_{k=1}^{\infty} a_k b_k$ converges. 
        \item   Give an example of two convergent series $\sum_{k=1}^{\infty} a_k$ and $\sum_{k=1}^{\infty} b_k$, such that $\sum_{k=1}^{\infty} a_k b_k$ diverges. 
        \item   Prove that if $\sum_{k=1}^{\infty} a_k$ and $\sum_{k=1}^{\infty} b_k$ converge absolutely, then $\sum_{k=1}^{\infty} a_k b_k$ converges absolutely. 
    \end{enumerate}

    \subsection{Solution (a)}
        \begin{gather}
            a_k = b_k = \frac{1}{k}\\
            a_k b_k = \frac{1}{k^2}
        \end{gather}

    \subsection{Solution (b)}
        \begin{gather}
            a_k = b_k = (-1)^{k + 1} \frac{1}{\sqrt{k}}\\
            a_k b_k = \frac{1}{k}
        \end{gather}

    \subsection{Solution (c)}
        \begin{proof}
            Note that these two converge absolutely.
            This means that their absolute value sums also converge.
            \begin{gather}
                \sum_{k = 1}^{\infty} \abs{a_k} \to L_1\\
                \sum_{k = 1}^{\infty} \abs{b_k} \to L_2
            \end{gather}

            Since they converge, we can use the contrapositive of the $k^{th}$-term test to determine what they converge to.
            \begin{gather}
                a_k \to 0; \abs{a_k} \to 0\\
                b_k \to 0; \abs{b_k} \to 0
            \end{gather}

            Take the member absolute value functions and multiply them together.
            \begin{equation}
                \abs{a_k} \abs{b_k} = \abs{a_k b_k} \to 0
            \end{equation}

            Since they converge to 0, there is a point after which $\abs{a_k} < 1$ and $\abs{b_k} < 1$.
            \begin{gather} \label{eq:4.48}
                \exists N \in \mathbb{N} : \forall k > N, \abs{a_k} < 1\\
                \exists N \in \mathbb{N} : \forall k > N, \abs{b_k} < 1
            \end{gather}

            We can multiply the last term of (\ref{eq:4.48}) by $\abs{b_k}$ on both sides.
            \begin{equation}
                \abs{b_k} > \abs{a_k} \abs{b_k} = \abs{a_k b_k}
            \end{equation}

            By the comparison test, this means that the sum of $(\abs{a_k b_k})$ over $k$ converges.
            From that, we can conclude that the sum of $(a_k b_k)$ over $k$ converges.
            \begin{gather}
                \sum_{k = 1}^{\infty} \abs{a_k b_k} \to L_3\\
                \sum_{k = 1}^{\infty} a_k b_k \to L_4
            \end{gather}
        \end{proof}
    \pagebreak

\section{Problem 4.22}
    Consider the sum $\sum_{k = 1}^{\infty} a_k$, and define $s_n$ to be  this series' $n^{th}$ partial sum; that is, $s_n = \sum_{k = 1}^{\infty} a_k$. 
    The series $\sum_{k=1}^{\infty} a_k$ is called Cesaro summable if the following converges:
    \begin{equation*}
        \lim_{n \to \infty} \frac{s_1 + s_2 + \dots + s_n}{n}
    \end{equation*}
    \begin{enumerate}[label=(\alph*)]
        \item   Prove that if $\sum_{k = 1}^{\infty} a_k$ converges, then this series is Cesaro summable.
        \item   Prove by example that if $\sum_{k = 1}^{\infty} a_k$ is Cesaro summable this does not imply that $\sum_{k = 1}^{\infty} a_k$ converges.
    \end{enumerate}

    \subsection{Solution (a)}
        \begin{proof}
            Since $\sum_{k = 1}^{\infty} a_k$ converges, its series of partial sums \seq{a} also converges.
            Recall Problem 3.29 in chapter 3, where we proved that if $\seq{a} \to a$, $b_n := \frac{a_1 + a_2 + \dots + a_n}{n} \to a$.
            Suppose we apply that to \seq{s}.
            \begin{gather}
                s_n \to s\\
                \frac{s_1 + s_2 + \dots + s_n}{n} \to s
            \end{gather}
            Conclsively, this means that \seq{s} converges.
        \end{proof}

    \subsection{Solution (b)}
        Take the case of $a_n = (-1)^n$.
        The sum diverges to a non-existant limit.
        However, dividing by the number of terms makes it converge to $0.5$.
    \pagebreak

\section{Problem 4.23}
    Show that if $\sum_{k = 1}^{\infty} a_k$ is conditionally convergent, then there exists a rearrangement of this sum which diverges to $\infty$. 

\section{Problem 4.24}
    Show that if $\sum_{k = 1}^{\infty} a_k$ is conditionally convergent, then there exists a rearrangement of this sum whose limit does not exist. 
    \pagebreak

\section{Problem 4.25}
    Assume that $\sum_{k = 1}^{\infty} a_k$ is conditionally convergent. 
    Define the limit superior of a sequence ($s_n$) (Notation: ``$\limsup_{n \to \infty} s_n$'') to be 
    \[
        \limsup_{n \to \infty} s_n = \lim_{n \to \infty} \left( \sup_{m \geq n} s_m \right)
    \]
    And define the limit inferior of a sequence ($s_n$) (Notation: ``$\liminf_{n \to \infty} s_n$'') to be 
    \[
        \liminf_{n \to \infty} s_n = \lim_{n \to \infty} \left( \inf_{m \geq n} s_m \right)
    \]
    Prove that for any $\alpha, \beta \in \mathbb{R} \cup \{\pm \infty\}$ with $\alpha \leq \beta$, there is a rearrangement of $\sum_{k = 1}^{\infty} a_k$ whose sequence of partial sums, $(s_n)$, have 
    \begin{align*}
        &\limsup_{n \to \infty} s_n = \beta &\andtxt &\liminf_{n \to \infty} s_n = \alpha
    \end{align*}

    \subsection{Solution}
        \begin{proof}
            We know that, for a conditionally convergent series, there exists a rearrangement that converges to any $L \in \mathbb{R} \cup \pm \infty$.
            Suppose defined a new variable $\delta$ such that $\delta = \frac{\abs{\alpha - \beta}}{4}$.
            We can separate $a_n$ into a sequence of its non-negative and negative terms, called \seq{p} and \seq{n} respectively.
            \begin{gather}
                \sum_{k=1}^{\infty} p_k \to \infty\\
                \sum_{k=1}^{\infty} n_k \to -\infty
            \end{gather}
            
            We can order \seq[k]{p} and \seq[k]{n} such that both converge to $0$.
            \begin{gather}
                \forall \varepsilon > 0, \exists N : \forall k > N, \abs{p_k} < \varepsilon\\
                \forall \varepsilon > 0, \exists N : \forall k > N, \abs{n_k} < \varepsilon
            \end{gather}

            There exists some $P_0$  and $N_0$ that we can sum up to the points $P_0$ and $N_0$ that $\abs{p_{P_0}} < \abs{\alpha - \beta}$ and $\abs{n_{N_0}} < \abs{\alpha - \beta}$. 
            \begin{equation}
                \sum_{k=1}^{P_0} p_k + \sum_{k=1}^{N_0} n_k = c_0
            \end{equation}

            Beyond that point, there exists some $P_1$ that we can sum to that will be greater than $\beta - \delta$.
            \begin{gather}
                \exists P_1 \geq 0 : c_0 + \sum_{k=P_0 + 1}^{P_1} p_k = b_1 \geq \beta - \delta
            \end{gather}

            From there, we can sum over \seq{n} until the total is less than $\alpha + \delta$.
            \begin{gather}
                \exists N_1 \geq 0 : \alpha < b_1 + \sum_{k=N_0 + 1}^{N_1} n_k = c_1 \leq \alpha - \delta
            \end{gather}

            We can continue this for any $m \geq 1$.
            \begin{gather}
                \forall m \in \mathbb{N}, \exists P_m \geq P_{m - 1} : \beta > c_{m - 1} + \sum_{k=P_{m-1}+1}^{P_m} p_k = b_m \geq b_{m - 1}\\
                \forall m \in \mathbb{N}, \exists N_m \geq N_{m - 1} : c_{m-1} > b_{m} + \sum_{k=N_{m-1}+1}^{N_m} n_k = c_m \leq \alpha
            \end{gather}

            Summing this up over all $m$, we end up with a sum that covers all of \seq{a} and which has a limit inferior of $\alpha$ and limit superior of $\beta$.
        \end{proof}
    \pagebreak

\section{Problem 4.26}
    Prove that if $\sum_{k = 1}^{\infty} a_k$ converges absolutely to $L$, then any rearrangement of this sum also converges to $L$. 

    \subsection{Solution}
        $\sum_{k = 1}^{\infty} a_k$ converges absolutely to $L$.
        This means that $\sum_{k = 1}^{\infty} \abs{a_k}$ converges to some $M \in \mathbb{R}$.
        \begin{equation}
            \sum_{k = 1}^{\infty} a_k \to L;
            \sum_{k = 1}^{\infty} \abs{a_k} \to M
        \end{equation}

\section{Problem 4.27}
    Prove that if each $a_k \geq 0$ and $\sum_{k = 1}^{\infty} a_k = \infty$, then any rearrangement of this sum also diverges to $\infty$. 
    \pagebreak

\section{Problem 4.29}
    There is a neat theorem which, in a small way, relies on the fact that the harmonic series $\sum_{k = 1}^{\infty} \frac{1}{k}$ diverges. 
    Here's the theorem: Every positive rational number can be written as a finite sum of distinct numbers of the form $1/n$. 
    For example, the positive rational number $\frac{7243}{4140}$ can be written like this: 
    \begin{equation*}
        \frac{7243}{4140} = \frac{1}{1} + \frac{1}{2} + \frac{1}{5} + \frac{1}{21} + \frac{1}{527} + \frac{1}{3,054,492}
    \end{equation*}
    Now, take a moment to appreciate how magical this seems, because once you see the algorithm which produces the decomposition, it might lose some of its shine. 
    Here's how I found the above decomposition of $\frac{7243}{4140}$ (which I basically chose at random): 
    \begin{gather*}
        \frac{7243}{4140} - 1 = \frac{3103}{4140}\\
        \frac{3103}{4140} - \frac{1}{2} = \frac{1033}{4140}\\
        \frac{1033}{4140} < \frac{1}{3} \andtxt \frac{1}{4}\\
        \frac{1033}{4140} - \frac{1}{5} = \frac{41}{828}\\
        \frac{41}{828} < \frac{1}{6}, \frac{1}{7}, \dots , \frac{1}{20}\\
        \frac{41}{828} - \frac{1}{21} = \frac{11}{5796}\\
        \frac{11}{5796} < \frac{1}{22}, \frac{1}{23}, \dots , \frac{1}{526}\\
        \frac{11}{5796} - \frac{1}{527} = \frac{1}{3,054,492}
    \end{gather*}
    \begin{enumerate}[label=(\alph*)]
        \item   If the fraction that we start with is much bigger than 1, like $\frac{244,406,536}{7}$, why is this not an issue? Why will subtracting off 1, and then 1/2, and then 1/3, and so on, eventually catch up to it? 
        \item   Prove that if this operation is performed on a fraction $\frac{p}{q}$ which is less than 1, then the numerator necessarily decreases (in the  above example, the numerators decreased at every stage: 7243,  3103, 1033, 41, 11, 1). Note: The algorithm always subtracts  off the biggest possible reciprocal. That is, if $\frac{p}{q} < 1$, then we subtract off $\frac{1}{n}$ where $\frac{1}{n} < \frac{p}{q} < \frac{1}{n - 1}$. 
        \item   Explain why this proves the theorem.
    \end{enumerate}

\end{document}